% !TEX encoding = UTF-8 Unicode
% !TEX spellcheck = en_US

\documentclass[a4paper, 11pt]{article}

\usepackage[utf8]{inputenc}
\usepackage[T1]{fontenc}
\usepackage{lmodern}
\usepackage[margin=2cm]{geometry}
\usepackage{enumitem}
\usepackage{fancyhdr}
\usepackage{lastpage}
\usepackage{amsmath}
\usepackage{graphicx}
\usepackage[hyphens]{url}
\usepackage{hyperref}



%% Specific commands
% the only place to specify version number (for consistency) - the date is the last modif for the code
\newcommand{\version}{1.10.00 (2023-03-17)}
\newcommand{\FullSWOF}{\emph{FullSWOF\_2D}}
\newcommand{\contactFullName}{Fr\'ed\'eric \textsc{Darboux} and Carine \textsc{Lucas}}
\newcommand{\contactEmail}{\href{mailto:Frederic.Darboux@inrae.fr}{frederic.darboux@inrae.fr} \href{mailto:carine.lucas@univ-orleans.fr}{carine.lucas@univ-orleans.fr}}
\newcommand{\MainWebSite}{\url{https://sourcesup.renater.fr/projects/fullswof-2d/}}
\newtheorem{cl}{Conclusion}

% Headers and footers
\lhead{\FullSWOF{} v\version}
\rhead{INRAE-IGE --- Univ. Orléans-IDP}
\cfoot{\thepage /\pageref{LastPage}}
\pagestyle{fancy}

% Lists: no spacing between items
\setlist{noitemsep}

\title{Improving the accuracy of floating-point computations
in~\FullSWOF{} - no. 2}

\author{\contactFullName\\ \contactEmail}
\date{2023-03-22} % Last modif for the note

\begin{document}

\maketitle

\thispagestyle{fancy}

\section*{Issue}

When running the 8~benchmarks of \FullSWOF{}, we obtain, with the \textgreater 2020 Apple chip~(ARM):
\begin{verbatim}
*** Comparison between reference and standard results ***
Emerged bump at rest: Maximum differences are: 9.000000e+00 (absolute) 
                                           and 5.786650e+00 (relative).
MacDonald: Rain fluvial Darcy-Weisbach: Results are identical.
MacDonald: Rain torrential Darcy-Weisbach: Results are identical.
MacDonald: Smooth transition with shock Manning: Results are identical.
MacDonald: Pseudo 2D torrential Manning: Results are identical.
MacDonald: Pseudo 2D fluvial Manning: Results are identical.
Dry dam break: Results are identical.
Thacker planar: Results are identical.
\end{verbatim}
See bug \href{https://sourcesup.renater.fr/tracker/?func=detail&aid=13511&group_id=895&atid=3706}{13511}. 

\section*{Testing \& Changes}
As for {\emph{FullSWOF\_1D}}, when looking at the results in details, we see that, 
for the bump at rest, the differences are very small, of order of $10^{-13} - 10^{-15}$~m!\\
$\rightarrow$ We performed a first change, in order to save the results up to $10^{-5}$~m, which is enough precise! 
This approximation is only performed in the output files, not for the computation. \\

The point is that the long doubles do not exist when using the new Apple configuration, they are replaced by doubles. \\
$\rightarrow$ We replaced long doubles by doubles in \texttt{hll.cpp}, \texttt{muscl.cpp}, \texttt{scheme.cpp}. \\

In the previous floating points application note (2014-02-05), we explained that long doubles were useful:
\begin{itemize}
\item to have the same results using Mac OSX, Linux and Windows, 
\item to have a symmetry in the dam break problem when water has an opposite velocity (from top-right to bottom-left). 
\end{itemize}
After performing the above modifications, we are able to check the symmetry of the results for the dam break problem.
The removal of the $\textgreater$~5~decimals widely helps the reproducibility of the results. \\
The results were validated with Mac OSX, Linux and Windows. 

\section*{Conclusions \& Recommendations}
We would like to draw the attention of users to these changes: they affect the results saved in the output files, as the are rounded to 5~decimals, 
but we believe that it will not affect the comparison to measurements, which are less precise. 

\end{document}

